% FILE: Medical terminology and definitions — glossary, definitions, medical terms
\chapter{Glossary of Medical and Scientific Terms}
\label{app:glossary}

This glossary defines medical, biochemical, immunological, and statistical terms used throughout this document. Terms are organized alphabetically. Where a term is used in a specialized sense specific to ME/CFS research, the ME/CFS-specific usage is indicated.

\section*{A}

\begin{description}
\item[Acetyl-CoA] Acetyl coenzyme~A; the two-carbon unit that enters the Krebs cycle, produced from pyruvate (via pyruvate dehydrogenase), fatty acids (via $\beta$-oxidation), or amino acids.
\item[ACTH] Adrenocorticotropic hormone; pituitary hormone that stimulates cortisol production by the adrenal cortex. Part of the HPA axis.
\item[Adenosine triphosphate (ATP)] The primary intracellular energy currency. Hydrolysis of ATP to ADP releases energy for cellular processes.
\item[Aerobic glycolysis] Glycolysis occurring in the presence of oxygen, with pyruvate entering the Krebs cycle rather than being converted to lactate. Distinct from anaerobic glycolysis.
\item[Allostatic load] The cumulative physiological cost of chronic stress and adaptation. In ME/CFS, elevated allostatic load reflects multi-system dysregulation.
\item[AMP-activated protein kinase (AMPK)] A cellular energy sensor activated by high AMP/ATP ratios. Triggers catabolic pathways to restore energy balance.
\item[Anaerobic threshold (AT)] The exercise intensity above which lactate production exceeds clearance, causing blood lactate to rise. Reduced in ME/CFS.
\item[Autoantibody] An antibody directed against the body's own proteins. Autoantibodies against adrenergic and muscarinic receptors have been identified in ME/CFS subsets.
\item[Autonomic nervous system (ANS)] The division of the nervous system controlling involuntary functions (heart rate, blood pressure, digestion) through sympathetic and parasympathetic branches.
\end{description}

\section*{B}

\begin{description}
\item[Baroreflex] The reflex arc that maintains blood pressure by adjusting heart rate and vascular tone in response to changes detected by arterial baroreceptors.
\item[Bell Disability Scale] A 0--100 rating scale for functional capacity in ME/CFS, where 100 represents normal function and 0 represents bedridden and dependent.
\item[Bistability] A dynamical system property in which two stable steady states coexist for the same parameter values. Transition between states requires a perturbation exceeding a threshold.
\item[Blood--brain barrier (BBB)] The selective permeability barrier formed by endothelial cells of cerebral capillaries, restricting passage of molecules and immune cells into the CNS.
\item[Butyrate] A short-chain fatty acid produced by gut bacteria from dietary fiber. Maintains intestinal barrier integrity and exerts anti-inflammatory effects.
\end{description}

\section*{C}

\begin{description}
\item[Cardiopulmonary exercise testing (CPET)] A diagnostic test measuring oxygen consumption, carbon dioxide production, and ventilation during graded exercise. The two-day CPET protocol is used to objectively demonstrate PEM.
\item[Cell danger response (CDR)] A conserved metabolic response to cellular threat, characterized by reduced mitochondrial function and extracellular ATP/ADP signaling. Proposed as a unifying mechanism in ME/CFS by Naviaux.
\item[Central sensitization] Enhanced responsiveness of central nervous system neurons to input, resulting in amplified pain and sensory perception.
\item[Cholinergic anti-inflammatory pathway] A neuroimmune reflex in which vagal efferent signals suppress macrophage TNF-$\alpha$ production through $\alpha$7 nicotinic acetylcholine receptors.
\item[Circadian rhythm] Endogenous biological oscillations with approximately 24-hour periodicity, governed by the suprachiasmatic nucleus and peripheral clock genes.
\item[Coenzyme Q10 (CoQ10)] Ubiquinone; an electron carrier in the mitochondrial electron transport chain between Complexes~I/II and Complex~III. Also functions as a lipid-soluble antioxidant.
\item[Corticotropin-releasing hormone (CRH)] Hypothalamic hormone that initiates the HPA axis cascade by stimulating pituitary ACTH release.
\item[Cortisol] The primary glucocorticoid hormone produced by the adrenal cortex. Regulates metabolism, immune function, and stress response. Mildly reduced in ME/CFS.
\item[Cytokine] A small signaling protein secreted by immune cells that mediates and regulates immunity, inflammation, and hematopoiesis. Key examples: IL-1$\beta$, IL-6, TNF-$\alpha$, IFN-$\gamma$, IL-10.
\end{description}

\section*{D--E}

\begin{description}
\item[DePaul Symptom Questionnaire (DSQ)] A validated self-report instrument for assessing ME/CFS symptom frequency and severity across multiple domains.
\item[Digital twin] A patient-specific computational model continuously updated with real-time data to predict individual disease trajectory and treatment response.
\item[Dysautonomia] Dysfunction of the autonomic nervous system. In ME/CFS, manifests as orthostatic intolerance, POTS, temperature dysregulation, and gastrointestinal dysmotility.
\item[Electron transport chain (ETC)] The series of protein complexes (I--IV) in the inner mitochondrial membrane that transfer electrons from NADH/FADH$_2$ to oxygen, generating the proton gradient that drives ATP synthesis.
\item[Energy envelope] The maximum amount of energy a patient can expend without triggering post-exertional malaise. Concept formalized by Jason et al.
\item[Epigenetics] Heritable changes in gene expression that do not involve changes to the DNA sequence. Includes DNA methylation, histone modification, and non-coding RNA regulation. Epigenetic alterations have been documented in ME/CFS.
\end{description}

\section*{F--G}

\begin{description}
\item[Fibromyalgia (FM)] A chronic pain condition frequently comorbid with ME/CFS, characterized by widespread musculoskeletal pain and tenderness.
\item[Fukuda criteria] A 1994 case definition for CFS requiring $\geq 6$ months of unexplained fatigue plus $\geq 4$ of 8 specified symptoms. Less specific than the ICC or CCC criteria.
\item[Glymphatic system] The brain's waste clearance system, dependent on aquaporin-4 water channels and active primarily during sleep. Impaired glymphatic function may contribute to unrefreshing sleep and cognitive dysfunction.
\item[Graded exercise therapy (GET)] A structured exercise program involving progressive increases in physical activity. Contraindicated in ME/CFS per NICE 2021 guidelines due to risk of harm.
\item[Gut--brain axis] The bidirectional communication system between the gastrointestinal tract and the central nervous system, involving neural (vagal), immune, endocrine, and microbial signaling pathways.
\end{description}

\section*{H--I}

\begin{description}
\item[Heart rate variability (HRV)] The variation in time intervals between consecutive heartbeats. Reduced HRV indicates autonomic dysfunction, particularly reduced parasympathetic tone. Consistently abnormal in ME/CFS.
\item[Hill function] A mathematical function describing cooperative binding or sigmoidal dose--response relationships: $f(x) = x^n / (K^n + x^n)$, where $n$ is the Hill coefficient and $K$ is the half-saturation constant.
\item[HPA axis] Hypothalamic--pituitary--adrenal axis; the neuroendocrine cascade (CRH $\to$ ACTH $\to$ cortisol) that regulates the stress response and cortisol production.
\item[Hypometabolism] A state of reduced metabolic activity. In ME/CFS, refers to the systemic downregulation of metabolic pathways documented by metabolomics studies.
\item[IDO (Indoleamine 2,3-dioxygenase)] An enzyme that converts tryptophan to kynurenine. Upregulated by IFN-$\gamma$, diverting tryptophan away from serotonin synthesis.
\item[Immune exhaustion] Progressive loss of effector function in T~cells or NK cells due to chronic antigen stimulation. Characterized by upregulation of inhibitory receptors (PD-1, Tim-3, LAG-3).
\item[International Consensus Criteria (ICC)] A 2011 case definition for ME requiring post-exertional neuroimmune exhaustion plus neurological, immune, and energy metabolism symptoms. More specific than the Fukuda criteria.
\end{description}

\section*{K--M}

\begin{description}
\item[Krebs cycle] (Also: citric acid cycle, tricarboxylic acid cycle.) The metabolic pathway in the mitochondrial matrix that oxidizes acetyl-CoA to CO$_2$, generating NADH, FADH$_2$, and GTP.
\item[Kynurenine pathway] The metabolic pathway converting tryptophan to nicotinamide (NAD$^+$ precursor) via kynurenine, 3-hydroxykynurenine, and quinolinic acid. Produces both neuroprotective and neurotoxic metabolites.
\item[Lipopolysaccharide (LPS)] A component of gram-negative bacterial cell walls. A potent activator of innate immunity through TLR4 signaling. Elevated circulating LPS (endotoxemia) may result from intestinal barrier dysfunction in ME/CFS.
\item[Low-dose naltrexone (LDN)] An opioid receptor antagonist used at sub-therapeutic doses (1--5\,mg) for immunomodulatory effects. Investigated as an off-label treatment for ME/CFS.
\item[Mast cell activation syndrome (MCAS)] A condition characterized by excessive mast cell mediator release, producing symptoms overlapping with ME/CFS (fatigue, brain fog, flushing, gastrointestinal disturbance).
\item[Metabolic trap] A hypothesis proposing that ME/CFS involves a stable biochemical state in which a metabolic pathway remains locked in a dysfunctional configuration due to bistability in enzyme kinetics.
\item[Michaelis--Menten kinetics] A model of enzyme-catalyzed reaction rates: $v = V_{\max}[S]/(K_m + [S])$, where $V_{\max}$ is the maximal rate and $K_m$ is the substrate concentration at half-maximal rate.
\item[Microglia] The resident immune cells of the central nervous system. Activated microglia produce neuroinflammatory mediators implicated in cognitive dysfunction and central sensitization.
\item[Mitochondria] Double-membrane organelles responsible for aerobic ATP production via the Krebs cycle and electron transport chain. Mitochondrial dysfunction is a central feature of ME/CFS pathophysiology.
\end{description}

\section*{N--O}

\begin{description}
\item[NAD$^+$/NADH] Nicotinamide adenine dinucleotide (oxidized/reduced forms). Essential electron carriers in cellular metabolism. The NAD$^+$/NADH ratio reflects cellular redox state and regulates metabolic flux.
\item[Natural killer (NK) cells] Innate immune lymphocytes that kill virally infected and transformed cells. Reduced NK cell cytotoxicity is one of the most replicated immunological findings in ME/CFS.
\item[Neuroinflammation] Inflammation within the central nervous system, mediated primarily by activated microglia and astrocytes. PET imaging has demonstrated neuroinflammation in ME/CFS patients.
\item[Orthostatic intolerance (OI)] Inability to tolerate upright posture, with symptoms (lightheadedness, palpitations, nausea, pre-syncope) that improve with recumbency. Affects the majority of ME/CFS patients.
\item[Oxidative phosphorylation (OXPHOS)] The process by which ATP is synthesized using energy from the mitochondrial proton gradient generated by the electron transport chain.
\item[Oxidative stress] An imbalance between reactive oxygen species production and antioxidant defense, resulting in cellular damage to lipids, proteins, and DNA.
\end{description}

\section*{P--R}

\begin{description}
\item[Pacing] An activity management strategy in which patients maintain energy expenditure within their energy envelope to prevent post-exertional malaise.
\item[Post-exertional malaise (PEM)] The hallmark symptom of ME/CFS: the persistence and accumulation of post-exertion byproducts --- metabolic waste, inflammatory mediators, oxidative damage --- that healthy physiology clears within hours but ME/CFS physiology cannot, producing symptom worsening typically delayed by 12--72 hours and lasting days to weeks.
\item[Postural orthostatic tachycardia syndrome (POTS)] A form of orthostatic intolerance defined by a sustained heart rate increase $\geq 30$\,bpm within 10 minutes of standing, without orthostatic hypotension.
\item[Proton-motive force] The electrochemical gradient across the inner mitochondrial membrane, composed of the membrane potential ($\Delta\Psi$) and the pH gradient ($\Delta$pH). Drives ATP synthesis by Complex~V.
\item[Reactive oxygen species (ROS)] Chemically reactive molecules derived from molecular oxygen, including superoxide ($\text{O}_2^{\bullet-}$), hydrogen peroxide (H$_2$O$_2$), and hydroxyl radical ($\bullet$OH). Generated as byproducts of mitochondrial electron transport.
\item[Rituximab] A monoclonal antibody targeting CD20 on B~cells, causing B~cell depletion. Investigated for ME/CFS based on an autoimmune hypothesis; phase~III trial did not demonstrate efficacy.
\end{description}

\section*{S--T}

\begin{description}
\item[Sensitivity analysis] A mathematical technique for determining which model parameters most strongly influence model outputs. Global sensitivity analysis (Sobol indices) partitions output variance across parameters and interactions.
\item[Separatrix] In dynamical systems theory, the boundary between basins of attraction. In the ME/CFS bistability model, the separatrix divides state space into regions that converge to the healthy state versus the disease state.
\item[SF-36 (Short Form Health Survey)] A 36-item questionnaire measuring eight domains of health-related quality of life. ME/CFS patients score lower than patients with most other chronic conditions.
\item[Short-chain fatty acids (SCFAs)] Fatty acids with fewer than 6 carbon atoms (acetate, propionate, butyrate) produced by bacterial fermentation of dietary fiber in the colon.
\item[Sickness behavior] A coordinated set of behavioral changes (fatigue, social withdrawal, anorexia, hyperalgesia) triggered by pro-inflammatory cytokines acting on the brain. Proposed as a model for ME/CFS symptomatology.
\item[Sympathetic nervous system (SNS)] The division of the autonomic nervous system responsible for ``fight-or-flight'' responses: increased heart rate, blood pressure, bronchodilation, and metabolic rate.
\item[Tilt table testing] A diagnostic procedure for orthostatic intolerance in which the patient is passively tilted from supine to upright while heart rate and blood pressure are monitored.
\item[Tryptophan] An essential amino acid and precursor for serotonin synthesis (via tryptophan hydroxylase) and kynurenine production (via IDO). Competition between these pathways is relevant to ME/CFS.
\end{description}

\section*{U--Z}

\begin{description}
\item[Ubiquinone] See Coenzyme Q10.
\item[Vagus nerve] The tenth cranial nerve; the primary parasympathetic nerve providing afferent and efferent innervation to thoracic and abdominal organs. Mediates the cholinergic anti-inflammatory pathway.
\item[Warburg effect] The metabolic shift from oxidative phosphorylation to aerobic glycolysis, originally described in cancer cells but also characteristic of activated immune cells.
\end{description}
